\documentclass[12pt]{article}
\usepackage[utf8]{inputenc}
\usepackage[T1]{fontenc}
\usepackage{amsmath, amssymb}
\usepackage{graphicx}
\usepackage{hyperref}
\usepackage{geometry}
\usepackage{listings}
\usepackage{xcolor}
\geometry{margin=1in}

\lstset{
    basicstyle=\ttfamily\small,
    breaklines=true,
    frame=single,
    backgroundcolor=\color{gray!10},
    keywordstyle=\color{blue},
    commentstyle=\color{gray},
    stringstyle=\color{teal}
}

\title{Web Technologies 1 (Frontend) \\ Week 4: CSS Layouts — Flexbox \& Grid}
\author{Dossymzhan}
\date{}

\begin{document}
\maketitle

\section{Overview}
Week 4 covers the \texttt{display} property, the difference between hiding elements, the two main modern layout systems — Flexbox (1D) and Grid (2D) — their terminology, key properties, and how to combine them.

\section{1. The \texttt{display} Property}
\texttt{display} controls how an element is shown on the page and how it interacts with its neighbors.
\begin{itemize}
    \item Default display: \texttt{block} or \texttt{inline}.
    \item The default can be changed, e.g. \texttt{display: inline} on \texttt{<li>} for a horizontal menu.
    \item Flexbox and Grid are turned on via \texttt{display: flex} and \texttt{display: grid}.
\end{itemize}
\textbf{Block-level elements:} start on a new line, take full width. Examples: \texttt{<div>}, \texttt{<h1>}–\texttt{<h6>}, \texttt{<p>}, \texttt{<form>}, \texttt{<header>}, \texttt{<footer>}, \texttt{<section>}.

\textbf{Inline elements:} stay in the same line, take only the needed width. Examples: \texttt{<span>}, \texttt{<a>}, \texttt{<img>}.

\section{2. \texttt{display: none} vs \texttt{visibility: hidden}}
\begin{itemize}
    \item \texttt{display: none} — element is removed; it takes up no space.
    \item \texttt{visibility: hidden} — element is hidden but its space is kept in the layout.
\end{itemize}

\section{3. Flexbox}
The CSS Flexible Box Layout Module (Flexbox) is a \textbf{one-dimensional} layout method for arranging items in a row or a column, without floats or manual positioning. Items grow to fill free space or shrink to fit the container.

With Flexbox you can:
\begin{itemize}
    \item center content horizontally and vertically;
    \item distribute space evenly between items;
    \item make columns the same height regardless of content;
    \item change the order/direction of items without editing the HTML.
\end{itemize}

\subsection{Flexbox Terminology}
\begin{itemize}
    \item \textbf{Flex container} — the parent with \texttt{display: flex}.
    \item \textbf{Flex items} — the direct children inside the container.
    \item \textbf{Main axis} — the primary axis along which flex items are laid out.
    \item \textbf{Cross axis} — perpendicular to the main axis.
    \item \textbf{main-start / main-end} — the starting/ending edges of the main axis.
    \item \textbf{main size} — total size of flex items along the main axis (width if row, height if column).
    \item \textbf{cross-start / cross-end} — starting/ending edges of the cross axis.
    \item \textbf{cross size} — total size of flex items along the cross axis.
\end{itemize}

\subsection{Key Flexbox Properties — Container}
\begin{center}
\begin{tabular}{ll}
\textbf{Property} & \textbf{Description} \\
\hline
\texttt{display} & Defines a flex container: \texttt{flex} or \texttt{inline-flex} \\
\texttt{flex-direction} & Main axis: \texttt{row}, \texttt{row-reverse}, \texttt{column}, \texttt{column-reverse} \\
\texttt{flex-wrap} & Wrapping: \texttt{nowrap}, \texttt{wrap}, \texttt{wrap-reverse} \\
\texttt{flex-flow} & Shorthand for \texttt{flex-direction} + \texttt{flex-wrap} \\
\texttt{justify-content} & Aligns items on the main axis (horizontally by default) \\
\texttt{align-items} & Aligns items on the cross axis (vertically by default) \\
\texttt{align-content} & Aligns flex lines when items wrap \\
\end{tabular}
\end{center}

\subsection{Key Flexbox Properties — Items}
\begin{center}
\begin{tabular}{ll}
\textbf{Property} & \textbf{Description} \\
\hline
\texttt{order} & Order in which items appear (default 0) \\
\texttt{flex-grow} & How much an item can grow relative to others (default 0) \\
\texttt{flex-shrink} & How much an item can shrink relative to others (default 1) \\
\texttt{flex-basis} & Initial size before growth/shrink (default \texttt{auto}) \\
\texttt{flex} & Shorthand for \texttt{flex-grow}, \texttt{flex-shrink}, \texttt{flex-basis} \\
\texttt{align-self} & Overrides \texttt{align-items} for a single item \\
\end{tabular}
\end{center}

\subsection{\texttt{flex-direction} and \texttt{justify-content}}
\begin{lstlisting}
.flex-container {
  display: flex;
  flex-direction: column; /* row (default) | row-reverse | column | column-reverse */
}
\end{lstlisting}
\texttt{justify-content} values: \texttt{flex-start} (default), \texttt{center}, \texttt{flex-end}, \texttt{space-between}, \texttt{space-around}, \texttt{space-evenly}.

\subsection{Flexbox Examples}
\begin{itemize}
    \item \textbf{Equal share:} \texttt{flex: 1} on every item — each takes an equal share of available space.
    \item \textbf{Auto margins:} \texttt{margin-left: auto;} pushes an item to the edge (e.g., a "Log in / Sign up" block pushed right in a navbar).
    \item \textbf{Column layouts:} \texttt{flex-direction: column;} stacks items vertically (e.g., a card with the button pinned to the bottom via \texttt{margin-top: auto;}).
    \item \textbf{Non-shrinking items:} \texttt{flex-shrink: 0;} keeps an item (e.g., an avatar) at a fixed size while sibling text wraps.
    \item \textbf{Wrapping:} \texttt{flex-wrap: wrap;} moves overflowing items to the next line.
    \item \textbf{Stretch alignment:} \texttt{align-items: stretch;} equalizes item heights.
\end{itemize}
\begin{lstlisting}
.card { flex-direction: column; }
.card .btn { margin-top: auto; }
.avatar { flex-shrink: 0; }
.buttons { margin-left: auto; }
\end{lstlisting}

\section{4. Grid Layout}
CSS Grid Layout is a \textbf{two-dimensional} grid-based system with rows \emph{and} columns, making it easier to build complex layouts without floats or positioning.

With CSS Grid you can:
\begin{itemize}
    \item define fixed, flexible, or proportional track sizes (\texttt{fr}, \texttt{auto}, \texttt{\%});
    \item place items into specific cells or span them across multiple rows/columns, regardless of source order;
    \item build complex layouts (dashboards, galleries, page templates) with consistent alignment.
\end{itemize}

\subsection{Grid Terminology}
\begin{itemize}
    \item \textbf{Grid Container} — parent with \texttt{display: grid}.
    \item \textbf{Grid Items} — direct children inside the grid.
    \item \textbf{Grid Track} — a row or a column.
    \item \textbf{Grid Area} — a rectangular space spanning multiple cells.
    \item \textbf{Grid Line} — dividing lines between cells.
    \item \textbf{Grid Cell} — a single "unit" of the grid.
    \item \textbf{Gap} — space between columns/rows.
\end{itemize}

\subsection{Special Units and Functions}
\begin{center}
\begin{tabular}{ll}
\textbf{Unit/Function} & \textbf{Description} \\
\hline
\texttt{fr} & Fraction of available space in the container (e.g. \texttt{1fr 2fr}) \\
\texttt{auto} & Size adjusts automatically to fit content \\
\texttt{min-content} & Track shrinks to the smallest size needed \\
\texttt{max-content} & Track expands to the largest size needed \\
\texttt{minmax(min, max)} & Track resizes between a min and max value \\
\texttt{fit-content(v)} & Track grows to fit content up to a maximum \\
\texttt{repeat(n, value)} & Repeats a track definition $n$ times (e.g. \texttt{repeat(3, 1fr)}) \\
\end{tabular}
\end{center}

\subsection{Key Grid Properties — Container}
\begin{center}
\begin{tabular}{ll}
\textbf{Property} & \textbf{Description} \\
\hline
\texttt{grid-template-columns} & Number and width of columns \\
\texttt{grid-template-rows} & Number and height of rows \\
\texttt{gap} / \texttt{row-gap} / \texttt{column-gap} & Spacing between tracks \\
\texttt{justify-items} / \texttt{align-items} & Aligns items horizontally/vertically within cells \\
\texttt{justify-content} & Aligns the entire grid along the row axis \\
\texttt{align-content} & Aligns the entire grid along the column axis \\
\texttt{grid-template-areas} & Defines the layout using named areas \\
\end{tabular}
\end{center}

\subsection{Key Grid Properties — Items}
\begin{center}
\begin{tabular}{ll}
\textbf{Property} & \textbf{Description} \\
\hline
\texttt{grid-column} & Column(s) an item should occupy \\
\texttt{grid-row} & Row(s) an item should occupy \\
\texttt{grid-area} & Item position by row/column, or a named area \\
\texttt{justify-self} & Aligns a single item horizontally in its cell \\
\texttt{align-self} & Aligns a single item vertically in its cell \\
\end{tabular}
\end{center}

\subsection{Spanning and Auto-fit}
\begin{lstlisting}
.grid {
  display: grid;
  grid-template-columns: repeat(auto-fit, minmax(110px, 1fr));
  gap: 10px;
}
.lead {
  grid-column: span 2; /* item takes 2 columns */
  grid-row: span 2;    /* item takes 2 rows    */
}
\end{lstlisting}
\texttt{grid-column: 1 / 4;} means "from line 1 to line 4". \texttt{repeat(auto-fit, minmax(110px, 1fr))} makes the number of columns follow the window width.

\subsection{Grid Template Areas}
\begin{lstlisting}
.page {
  display: grid;
  grid-template-columns: 120px 1fr;
  grid-template-rows: 50px 1fr 40px;
  grid-template-areas:
    "header header"
    "sidebar main"
    "footer footer";
  gap: 10px;
}
.header  { grid-area: header;  }
.sidebar { grid-area: sidebar; }
.main    { grid-area: main;    }
.footer  { grid-area: footer;  }
\end{lstlisting}
This names each grid area directly in the CSS layout string, giving a readable "map" of the page (header / sidebar + main / footer).

\section{5. Flexbox vs Grid}
\begin{center}
\begin{tabular}{lll}
\textbf{Feature} & \textbf{Flexbox (1D)} & \textbf{Grid (2D)} \\
\hline
Layout direction & Row \emph{or} Column & Rows \emph{and} Columns \\
Use case & Navbars, simple alignment & Complex page layouts \\
Approach & Content-first & Layout-first \\
Responsiveness & Easy with \texttt{flex-wrap} & Easy with \texttt{auto-fit}/\texttt{auto-fill} \\
\end{tabular}
\end{center}

\subsection{When to Use Which}
\begin{itemize}
    \item \textbf{Use Flexbox for:} navbars, menus, buttons, form controls, one-direction layouts.
    \item \textbf{Use Grid for:} full page layout, dashboards, galleries, complex responsive designs.
\end{itemize}

\section{6. Combining Flexbox and Grid}
They work well together:
\begin{itemize}
    \item Grid for the large-scale page layout; Flexbox for aligning content inside individual grid cells.
    \item A grid item can itself be a flex container.
    \item A flex item can itself be a grid container.
\end{itemize}
\begin{lstlisting}
/* page layout: Grid */
.page {
  display: grid;
  grid-template-columns: 120px 1fr;
  grid-template-rows: 60px 1fr 40px;
  gap: 10px;
}

/* inside cells: Flexbox */
.header {
  display: flex;
  justify-content: space-between;
  align-items: center;
}
.main {
  display: flex;
  gap: 10px;
}
.card { flex: 1; }
\end{lstlisting}

\section{Summary}
Week 4 covers the \texttt{display} property and how it differs from visibility, then the two modern CSS layout systems: one-dimensional \textbf{Flexbox} (container/item terminology, main/cross axis, \texttt{flex-direction}, \texttt{justify-content}, growing/shrinking items) and two-dimensional \textbf{Grid} (tracks, areas, lines, cells, \texttt{fr} units, \texttt{repeat()}/\texttt{minmax()}, \texttt{grid-template-areas}), when to use each, and how to combine them for real page layouts.

\end{document}
